\documentclass{situationist-preprint}

\title{The Explanatory Gap as a Gap Between Modes of Presentation}
\msid{MND-01-b}
\msfamily{MND-01}
\msdiscipline{Philosophy of Mind}
\mstype{Companion manuscript}
\mskeywords{explanatory gap; modes of presentation; acquaintance; Russellian monism; phenomenal concepts; Ramsey sentences; meta-problem}

\begin{document}
\maketitle

\begin{abstract}
MND-01 holds that consciousness is the intrinsic nature of physical structure and that the explanatory gap arises from two modes of presentation of the same facts. We develop the second half of the thesis. Physical truths present their subject matter structurally, through the roles that physical terms describe; experience presents it intrinsically, by acquaintance. We show that the a priori non-derivability of phenomenal truths from physical truths follows from this difference alone, that it persists under any structural extension of physics, and that it reduces to a single kind of a posteriori identification: that the occupant of a role is the nature presented in acquaintance. Unlike the identities of type-B physicalism, this identification does not pair two rival characterizations of one property, and so leaves no residual question of why a given structure should feel as it does. We apply the account to Mary, zombies, the inverted spectrum and the meta-problem.
\end{abstract}

\section{Introduction}\label{s:intro}

\citet{Levine1983} observed that even a true psychophysical identity, such as the identity of pain with the firing of C-fibres, would leave something unexplained. The identity of water with H$_2$O explains why water boils at a certain temperature and dissolves salt, because once we know what H$_2$O does we can see that anything that does it is water. Nothing we could learn about what C-fibre firing does would let us see that it must feel painful. There is, Levine concluded, an explanatory gap between the physical and the phenomenal, and the gap persists whether or not the identity is true.

The gap admits two readings. On the \emph{ontological} reading, it reflects a gap in the world: phenomenal facts are further facts, not fixed by the physical ones, and the failure of derivation is the epistemic shadow of that metaphysical independence \citep{Chalmers1996}. On the \emph{epistemic} reading, the gap reflects a difference in how the same facts are known: there is one subject matter and two ways of thinking about it, and the ways are not a priori connected \citep{Loar1990,Papineau2002}. MND-01 takes the epistemic reading, but with a content that differs from the standard version of it. The principal manuscript states the thesis formally: given materialism, the zombie hypothesis $P \wedge \neg Q$ is primarily possible but not secondarily possible (MND-01-a, Corollary~4.4), and it attributes the difference to the fact that the terms of $P$ pick out their referents by structural roles that the intrinsic natures presented in experience actually occupy. That corollary says \emph{where} the gap lies; the present manuscript says \emph{what} it is, and why the difference between the two modes of presentation is sufficient to produce it.

The thesis we defend is this. Physical truths present their subject matter in a \emph{structural} mode: they say what plays which role in a network of lawful relations. Phenomenal truths, as known to the subject, present their subject matter in an \emph{intrinsic} mode: by acquaintance with how something is in itself. The explanatory gap is the gap between these two modes. It is real, ineliminable by further physics, and wholly epistemic. It consists in the a priori independence of two modes of presentation that, given the Russellian monism of MND-01, present the same portion of reality.

The thesis has a rival within the epistemic camp, the phenomenal concept strategy, and part of our task is to show that the two are different and that the difference matters. On the phenomenal concept strategy, physical concepts and phenomenal concepts pick out the same physical or functional property in two ways, and the gap is explained by the peculiarities of phenomenal concepts. On the present view, the physical concepts do not characterize the property that the phenomenal concepts present; they characterize only the role it occupies. This removes the feature of type-B identities that Levine found unsatisfying, and it escapes the main argument against the phenomenal concept strategy.

Section~\ref{s:history} reviews the debate over the gap. Section~\ref{s:frame} states the framework. Section~\ref{s:arg} proves that the gap follows from the difference in modes, that it persists, and that it reduces to one kind of a posteriori identification. Section~\ref{s:cases} applies the account to four cases. Section~\ref{s:obj} considers objections, and Section~\ref{s:not} says what the account does not establish.

\section{The debate over the gap}\label{s:history}

\subsection{Levine and the demand for derivation}

Levine's argument rests on a model of reductive explanation. A higher-level phenomenon is explained by a lower-level one when the lower-level description, together with the concept of the higher-level phenomenon, lets us derive that the phenomenon occurs. The concept of water is, at least in part, the concept of whatever is the clear, drinkable liquid that fills the lakes; the chemistry of H$_2$O shows that H$_2$O is such a liquid; so water is H$_2$O, and the identity is explained rather than merely asserted \citep{Levine1993}. The concept of pain is not, or not only, the concept of whatever plays a causal role. It includes a conception of what pain feels like, and nothing in the neurophysiology of C-fibres lets us derive that their firing feels like that. Levine was careful to distinguish this epistemic conclusion from the metaphysical conclusion that pain is not C-fibre firing, and in later work he argued that the gap is evidence of something unresolved in our conception of the mind without being evidence of dualism \citep{Levine2001}.

\subsection{Two-dimensionalism and the rejection of the derivation model}

The model of reductive explanation that Levine relied on was made precise by \citet{ChalmersJackson2001}. They argued that every macroscopic truth that is not phenomenal or indexical is a priori entailed by the conjunction of microphysical truths, phenomenal truths, an indexical locator and a \enquote{that's all} clause, and that reductive explanation consists in such entailment. Phenomenal truths are the exception, and the exception, they argued, shows that phenomenal truths are not reductively explicable and, given the link between a priori entailment and necessitation, that materialism is false.

\citet{BlockStalnaker1999} rejected the model. They argued that ordinary reductive explanations do not proceed by a priori entailment from microphysics: we do not have, and do not need, conceptual analyses of \enquote{water} that would let us deduce its presence from a microphysical description. Identities in science are accepted on the ground that they simplify and unify our picture of the world, and an identity of pain with a neural state could be accepted on the same ground. If this is right, there is no special gap for consciousness: every identity is a posteriori and every explanation leaves something to be accepted on grounds of simplicity. The reply of \citet{ChalmersJackson2001} was that the derivation is available in principle for water, given enough information about the role, whereas it is unavailable in principle for pain, and that the difference is what needs explaining. The present account accepts this reply. The gap is special, and an account of it must say why.

\subsection{The phenomenal concept strategy}

The most developed epistemic account is the phenomenal concept strategy. \citet{Loar1990} proposed that phenomenal concepts are recognitional concepts that refer directly to physical-functional properties without presenting them under a physical or functional description. Since phenomenal and physical concepts are conceptually isolated, no a priori connection holds between them, and the gap follows; but the concepts may co-refer, and so materialism is untouched. \citet{Papineau2002} developed the view with a quotational model of phenomenal concepts, on which a phenomenal concept embeds a version of the experience it refers to, and \citet{Balog2012} with an account based on acquaintance. The collection adopts a version of this account for Mary and for zombies: Mary acquires on her release a phenomenal concept of a fact she already knew (MND-04-a, Theorem~3.3), and phenomenal concepts refer directly, without a descriptive mode of presentation that would connect them a priori with physical concepts (MND-05-a, Definition~2.1 and Assumption~2.2).

The strategy has two difficulties that bear on our argument. The first is Levine's residual question. If a phenomenal concept and a neural concept pick out one property, the identity is a posteriori, and its a posteriori status is explained by the isolation of the concepts. But the neural concept does characterize the property: it says that the property is a certain pattern of activity with a certain functional profile. And the phenomenal concept characterizes it too: it presents it as feeling a certain way. We are left with one property characterized in two ways, and the question why something with the first character should have the second remains \citep{Levine2001}. The second difficulty is the master argument of \citet{Chalmers2007}. Either the features of phenomenal concepts that are supposed to explain the gap are themselves physically explicable, or they are not. If they are, a zombie would have concepts with the same features, and so the features cannot explain our epistemic situation with respect to consciousness, which the zombie does not share. If they are not, the strategy has relocated the gap rather than closed it.

\citet{Tye2009}, who had defended the strategy, abandoned it, arguing that phenomenal concepts as the strategy requires them do not exist. \citet{Stoljar2005} argued that the strategy, even if successful, explains less than it seems to, because the relevant features of phenomenal concepts can be possessed only if the concepts are already of phenomenal properties, which is what was to be explained.

\subsection{The Russellian turn}

A different response begins from the observation of \citet{Russell1927} that physics describes the world structurally. If it does, the physical truths in the antecedent of Levine's conditional leave out the intrinsic natures that realize physical structure, and the failure of derivation may be due to that omission rather than to the existence of non-physical facts. \citet{Stoljar2006} developed a version of this response as the \emph{ignorance hypothesis}: we are ignorant of a type of physical truth, and if we knew it, the gap would close. The ignorance hypothesis is non-committal about what the missing truths are. Russellian monists are committed: the missing truths concern intrinsic natures that are phenomenal or protophenomenal \citep{Chalmers1996,Goff2017}. MND-01 adopts the protophenomenal version, on which the intrinsic natures of fundamental entities are not experiences but ground experiences in suitable arrangement (MND-18-a, Corollary~3.3; MND-01-d, Proposition~1.1; Volume~III of MND-01, Definition~21.2).

The Russellian turn changes the question about the gap. On the phenomenal concept strategy, the gap is between two concepts of one structurally characterized property. On the Russellian view, the gap is between a concept that characterizes a role and a concept that presents what occupies it. The rest of this manuscript develops that difference.

\subsection{The meta-problem}

Finally, \citet{Chalmers2018} argued that any theory of consciousness should explain why we judge that there is a hard problem: why we find it puzzling that physical processes are accompanied by experience. The judgment is a physical event, and its causes are in principle discoverable by cognitive science. Illusionists take the meta-problem to be the only real problem \citep{Frankish2016}. MND-10-a argues that illusionism is false, since the seeming it posits is either itself phenomenal or fails to explain the datum (MND-10-a, Lemma~4.2). We do not rely on that result here, but an account of the gap should solve the meta-problem, and Section~\ref{s:cases} argues that the present one does.

\section{Framework}\label{s:frame}

We use the notation of MND-01-a. $P$ is the conjunction of microphysical truths, and $Q$ an arbitrary phenomenal truth. A truth is \emph{a priori} for a subject if it can be known with justification independent of experience, given possession of the concepts in which it is stated; \enquote{experience} here means evidence, not the phenomenal states that are our topic.

\begin{definition}[Mode of presentation]\label{d:mode}
A \emph{mode of presentation} is a way in which an object, property or fact is given to a thinker in thought. Two thoughts that concern the same object, property or fact present it in different modes if and only if a thinker who possesses both and reflects on them can rationally accept one and doubt or reject the other.
\end{definition}

The criterion is Frege's \citep{Frege1892}: \enquote{Hesperus is Hesperus} and \enquote{Hesperus is Phosphorus} concern one planet, and a rational thinker can accept the first and doubt the second, so \enquote{Hesperus} and \enquote{Phosphorus} present the planet in different modes. Definition~\ref{d:mode} does not commit us to any theory of what modes of presentation are; it fixes them by their role in rational cognition. In particular, it allows a mode of presentation to be non-descriptive. The phenomenal concepts of MND-05-a refer directly, without a \emph{descriptive} mode of presentation, but by Definition~\ref{d:mode} they still present their referents in a mode, since a subject can rationally accept a thought containing one of them and doubt a thought about the same property that contains a physical concept.

\begin{definition}[Structural mode]\label{d:struct}
A thought presents its subject matter in the \emph{structural mode} if its content is fixed by a Ramsey description: it concerns whatever entities and properties stand in a specified pattern of nomic, causal, spatiotemporal and mathematical relations, and it says nothing about those entities and properties beyond their places in the pattern.
\end{definition}

The theoretical terms of physics are defined, on the account of \citet{Lewis1970}, by their places in the Ramsey sentence of the theory: \enquote{charge} denotes whatever property occupies the charge-role in the pattern of relations that electromagnetism describes. A thought stated in such terms concerns its subject matter only as the occupant of a role. This is the content of Structuralism about physics, \ref{p:struct-a} below, and of P2 of MND-01-a.

\begin{definition}[Intrinsic mode]\label{d:intr}
A thought presents its subject matter in the \emph{intrinsic mode} if its content is fixed by the subject's acquaintance with an instance of the property it concerns, and it concerns the property as it is in itself and not as the occupant of a role.
\end{definition}

Acquaintance in the sense of \citet{Russell1911} is the most direct cognitive relation in which a subject can stand to an object, a relation in which the object is presented without the mediation of a description. Russell held that we are acquainted with sense-data and universals; we hold only that a subject is acquainted with the phenomenal character of her current, attended experiences, as MND-10-a's principle of introspective reliability presupposes (MND-10-a, P1). Definition~\ref{d:intr} does not require that acquaintance reveal everything about the property: we return to this in Section~\ref{s:obj}.

We now state three principles.

\begin{principle}[Structuralism about physics]\label{p:struct-a}
The microphysical truths $P$ present their subject matter in the structural mode.
\end{principle}

\ref{p:struct-a} is P2 of MND-01-a restated in terms of modes of presentation. It says nothing about what realizes the structure, and so does not say that the realizers are unknown in every sense: only that $P$ does not present them otherwise than as role-occupants.

\begin{principle}[Structural closure]\label{p:closure}
Whatever is a priori derivable from truths stated in the structural mode, together with indexical and \enquote{that's all} information, is itself stated, or statable, in the structural mode.
\end{principle}

\ref{p:closure} is the epistemic counterpart of the observation, made in MND-01-a, that no derivation from structure alone yields a truth that is not about structure. Its justification is that a priori reasoning is closed under the operations by which content is fixed: an a priori consequence of a set of premises concerns nothing that the premises do not concern, except through the logical and mathematical vocabulary that a priori reasoning adds, and that vocabulary is itself structural. \citet{ChalmersJackson2001} accept the principle in effect when they argue that macroscopic truths are derivable from $P$ because the concepts of macroscopic kinds are concepts of roles.

\begin{principle}[Non-structural content of acquaintance]\label{p:acq}
A phenomenal truth $Q$, as known to a subject acquainted with the experience it concerns, presents that experience in the intrinsic mode, and its content is not statable in the structural mode.
\end{principle}

\ref{p:acq} is supported by MND-10. A complete relational description of a quality space determines the assignment of qualia to places only up to automorphism and never determines which intrinsic character the occupant of a place has (MND-10-a, Proposition~6.1); qualia are therefore not constituted by their places in a quality space (MND-10-a, Lemma~5.2). A subject who knows that her experience is like \emph{this} knows something about it that no structural description, however complete, states.

Finally we take the metaphysical half of MND-01 as given.

\begin{assumption}[Russellian monism]\label{a:rm}
The physical terms in $P$ refer to whatever actually occupies the roles that $P$ describes. The occupants of those roles have intrinsic natures that are protophenomenal, and every phenomenal property is wholly grounded in these natures together with the structure in which they stand.
\end{assumption}

Assumption~\ref{a:rm} is Corollary~4.3 of MND-01-a, which holds given materialism, together with the protophenomenal reading of the view required by MND-18 (MND-18-a, Corollary~3.3; MND-18-b, Definition~1.1). Materialism here is the physicalism of MND-02, which counts categorical realizers of dynamical roles as physical (MND-02-a, Definition~2.2(b)) and is true given its premises (MND-02-a, Theorem~3.6).

We say that two truths \emph{present the same portion of reality} if what the second concerns is wholly grounded in what the first concerns, so that the obtaining of the first's subject matter, as it actually is, necessitates the obtaining of the second's and adds nothing further to what there is. This is the sense in which MND-01 speaks of the same facts: not that $P$ and $Q$ express one fine-grained proposition, which they plainly do not, but that the reality $Q$ describes is nothing over and above the reality $P$ describes, taken with its actual intrinsic character.

\section{The argument}\label{s:arg}

\subsection{The gap}

\begin{lemma}\label{l:gap}
The conditional $P \to Q$ is not a priori.
\end{lemma}

\begin{proof}
Suppose it were. Then $Q$ would be a priori derivable from $P$ (with indexical and \enquote{that's all} information, if needed). By \ref{p:struct-a}, $P$ is stated in the structural mode, so by \ref{p:closure}, $Q$ would be statable in the structural mode. By \ref{p:acq}, it is not.
\end{proof}

Lemma~\ref{l:gap} is P1 of MND-01-a, and so the gap itself. The point of deriving it is to show what it rests on. The proof uses no premise about the metaphysics of consciousness: it uses only facts about the modes in which $P$ and $Q$ present their subject matter. A dualist and a Russellian monist can both accept it. What distinguishes them is what they say about the subject matter, and to that we now turn.

\begin{lemma}\label{l:same}
Given Assumption~\ref{a:rm}, $P$ and $Q$ present the same portion of reality: $P$, with its terms taken to refer to what actually occupies the roles it describes, necessitates $Q$, and what $Q$ concerns is wholly grounded in what $P$ concerns.
\end{lemma}

\begin{proof}
By Assumption~\ref{a:rm}, the terms in $P$ refer to the actual occupants of physical roles, and those occupants have protophenomenal natures in which, together with structure, every phenomenal property is wholly grounded. So the subject matter of $P$, as it actually is, includes the full ground of the phenomenal property that $Q$ attributes, and since full grounding entails necessitation, every world at which that subject matter obtains is a world at which $Q$ is true. Since what is wholly grounded adds nothing to its ground, $Q$ concerns nothing over and above what $P$ concerns.
\end{proof}

\begin{theorem}\label{t:modes}
Given Assumption~\ref{a:rm} and \ref{p:struct-a}--\ref{p:acq}, the conditional $P \to Q$ is necessary and not a priori, and its failure to be a priori is fully explained by the difference between the structural mode in which $P$ presents its subject matter and the intrinsic mode in which $Q$ presents it. There are not two portions of reality corresponding to $P$ and to $Q$.
\end{theorem}

\begin{proof}
Necessity is Lemma~\ref{l:same}; the failure of apriority is Lemma~\ref{l:gap}. The proof of Lemma~\ref{l:gap} used only \ref{p:struct-a}--\ref{p:acq}, which concern modes of presentation, so the failure of apriority is explained by the difference in modes and requires no difference in subject matter. By Lemma~\ref{l:same} there is none.
\end{proof}

Theorem~\ref{t:modes} is the thesis of MND-01 in the form proper to the present manuscript. It agrees with MND-01-a, Corollary~4.4: since $P \to Q$ is necessary, $P \wedge \neg Q$ is not secondarily possible; since it is not a priori, $P \wedge \neg Q$ is not ruled out a priori, and on the two-dimensional framework this is what its primary possibility consists in. It also agrees with MND-05: the zombie hypothesis is ideally conceivable but not metaphysically possible (MND-05-a, Theorem~3.3), and the failure of the conceivability--possibility principle is confined to propositions formulated with phenomenal concepts (MND-05-a, Corollary~3.4). Theorem~\ref{t:modes} adds an account of why those concepts are special: they are the concepts in which reality is presented in the intrinsic mode.

\subsection{Persistence}

A natural hope is that the gap will close as physics progresses. The hope is misplaced, for a reason that Theorem~\ref{t:modes} makes precise.

\begin{proposition}[Persistence]\label{p:persist}
Let $P^{+}$ be the conjunction of truths of any extension of physics that states its truths in the structural mode. Then $P^{+} \to Q$ is not a priori.
\end{proposition}

\begin{proof}
The proof of Lemma~\ref{l:gap} applies with $P^{+}$ for $P$, since it used of $P$ only that it is stated in the structural mode.
\end{proof}

Proposition~\ref{p:persist} distinguishes the present account from the ignorance hypothesis of \citet{Stoljar2006}, on which the gap would close if we came to know a currently unknown type of physical truth. On the present account, there is indeed a type of physical truth that physics does not state, namely truths about the intrinsic natures of role-occupants. But it is not a type of truth that a further physics, stated as physics is stated, could supply, since any theory that characterizes its posits by their nomic roles presents them in the structural mode. The missing truths are known, in so far as they are known at all, through acquaintance, and only for the arrangements of natures that ground our own experiences. This is the sense in which the gap is ineliminable without being ontological.

The point bears on the definition of the physical. MND-02 defines physical properties by reference to current physics and its foreseeable extensions and counts categorical realizers of dynamical roles among the fundamental properties of such an extension (MND-02-a, Definition~2.2(b)). A foreseeable extension may thus include truths that its realizers exist and occupy roles; Proposition~\ref{p:persist} says that its statements about them will still be in the structural mode. The physical facts include the intrinsic natures; the physical \emph{truths}, as physics states them, present them only as occupants.

\subsection{Reduction to one identification}

Levine's complaint against type-B physicalism was that its identities are brute: they assert that a property characterized physically is a property characterized phenomenally, without letting us see why. We now show that, on the present account, the a posteriori element in the relation between $P$ and $Q$ takes a specific form, and that the form does not invite the complaint.

\begin{definition}[Occupant identification]\label{d:occ}
An \emph{occupant identification} is a truth of the form \enquote{the occupant of role $R$ is \emph{this}}, where $R$ is specified in the structural mode and \enquote{\emph{this}} is a concept of the intrinsic mode.
\end{definition}

\begin{proposition}\label{p:occ}
Given Assumption~\ref{a:rm} and \ref{p:struct-a}--\ref{p:acq}, there is a conjunction $M$ of occupant identifications such that $(P \wedge M) \to Q$ is a priori, while $M$ is not a priori.
\end{proposition}

\begin{proof}
By Assumption~\ref{a:rm}, the phenomenal property that $Q$ attributes is wholly grounded in the protophenomenal natures and structure of some arrangement of role-occupants, and in the actual world the arrangement that grounds it is the arrangement with which the subject is acquainted, under its intrinsic aspect, when she has the experience. Let $R$ be the structural role of that arrangement as specified in $P$, and let $M$ be the identification \enquote{the occupant of $R$ is \emph{this}}, with \enquote{\emph{this}} the intrinsic-mode concept by which the subject thinks of her experience. Given $P$, the occupant of $R$ is instantiated; given $M$, it is \emph{this}; that \emph{this} is instantiated is $Q$, or entails it by the content of the concept. So $(P \wedge M) \to Q$ is a priori. $M$ is not a priori, since by Lemma~\ref{l:gap} $P \to Q$ is not, and if $M$ were a priori so would be $P \to Q$.
\end{proof}

Proposition~\ref{p:occ} says that the whole a posteriori residue in the passage from physics to experience consists of occupant identifications. These are of a familiar kind. \enquote{The occupant of the role of causing these symptoms is \emph{this} virus}, said by someone who has isolated it, is one; so is \enquote{the occupant of the evening-star role is \emph{this} planet}, said by an astronomer looking through a telescope. In each, one term specifies a role and the other presents an occupant by direct cognitive contact. What is distinctive in the phenomenal case is only that the direct contact is acquaintance with the occupant's intrinsic character rather than perception of its relational properties.

The difference from type-B identities is now visible. A type-B identity says that a property characterized by its neural and functional profile is a property characterized by its phenomenal feel. Both sides characterize the property, and the question why something characterized in the first way should be characterized in the second arises because both characterizations purport to say what the property is like. An occupant identification says that the occupant of a role is a property characterized by its phenomenal feel. The first side does not say what the property is like at all; it says only what it does. There is no rival characterization with which the phenomenal one could fail to cohere. To ask why the occupant of role $R$ should be like \emph{this} is to ask why the realizer of a structure has the intrinsic nature it has, and that is a question that arises for every realizer of every structure, not a special question about consciousness. It is answered, so far as it can be answered, by the grounding principles that pair arrangements with the qualities they ground, which are the subject of Volume~IV of MND-01 and of the palette problem discussed in MND-01-d.

\begin{corollary}\label{c:noresid}
On the present account there is no residual explanatory question about consciousness beyond the question of the grounding principles: given $P$, the occupant identifications $M$ and the grounding principles, $Q$ is explained.
\end{corollary}

\begin{proof}
By Proposition~\ref{p:occ}, $Q$ follows a priori from $P$ and $M$. $M$ is a posteriori but is not brute in Levine's sense, since its structural side does not characterize the occupant; the only further question it raises, why this arrangement grounds this quality, is a question about the grounding principles.
\end{proof}

Corollary~\ref{c:noresid} does not say that we now possess the grounding principles. MND-01-d argues that the boundary and palette problems are problems of stating them, not of showing that they can exist (MND-01-d, Proposition~2.1), and Volume~IV of MND-01 takes up that task, including the principle that maximal integrated arrangements fix the boundaries of subjects. Our claim is that once those principles are stated, nothing about consciousness remains that is not explained, and that the explanatory gap, which will remain, is not a sign that something does.

\subsection{Asymmetry}

The gap runs in one direction more fully than in the other. From $P$ one cannot derive $Q$. Can one derive from $Q$, as known by acquaintance, the structural truths about the experience?

\begin{proposition}\label{p:asym}
A subject acquainted with her experiences can know a priori, from her acquaintance with them, some structural truths about them, namely the relations among their qualities that a relational description of her quality spaces states; she cannot know from acquaintance alone which physical role the arrangement grounding them occupies.
\end{proposition}

\begin{proof}
For the first part: by MND-10-a, Theorem~5.4, the structure of each modality's qualia can be fully described relationally, and the relations among qualities, such as relative similarity and betweenness, are among the things a subject attending to her experiences can judge reliably (MND-10-a, P1). For the second part: by \ref{p:struct-a}, physical roles are specified in the structural mode by their places in the nomic network of physics, and nothing in the intrinsic character of an experience fixes which nomic role its ground occupies; the identification of that role requires the a posteriori occupant identification of Proposition~\ref{p:occ}.
\end{proof}

The asymmetry explains a familiar feature of the debate. The structure of phenomenal space, the similarity and difference relations among colors, tones and tastes, is accessible to psychophysics and has been mapped in detail; it is a structural matter and so is shared between the two modes. The intrinsic character of each quality is not shared. Hence the partial success of structural specification in the case of the bat, which MND-03 develops: much of the character of echolocation can be specified structurally, as a sonar-structured spatial field, and what cannot be specified is the intrinsic character of the qualities that occupy the places of that structure (MND-03-a, Theorem~5.3).

\section{Cases}\label{s:cases}

\subsection{Mary}

Mary knows $P$, and every structural truth about color vision, before her release. MND-04 argues that she learns no new fact when she sees red; she acquires a recognitional ability and with it a phenomenal concept of a fact she already knew (MND-04-a, Theorem~3.3). The present account supplies the reason. Before her release Mary knows the fact that the arrangement grounding her red-experiences is instantiated, presented in the structural mode, as the occupant of a role. After her release she is acquainted with it, and so can think of it in the intrinsic mode. The fact is the same, by Lemma~\ref{l:same}; the mode is new. Her new phenomenal concept co-refers with a physical concept she already had, and the identity of the two is a posteriori (MND-04-b, Propositions~1.1 and~1.2); by Proposition~\ref{p:occ}, the identity is an occupant identification. Her continuing sense that the physical facts leave something out (MND-04-b, Corollary~2.1) is the gap of Lemma~\ref{l:gap}, which persists after her release because her new concept does not convert the structural mode of $P$ into the intrinsic mode.

The account also explains why Mary seems to learn something, and why the appearance is not wholly misleading. She acquires knowledge of the intrinsic character of an occupant that she knew only as an occupant. On the coarse individuation of facts used throughout MND-01, this is not knowledge of a new fact. On a finer individuation, it is new knowledge, and it is the only new knowledge that her situation makes available. What she learns is what the occupant of the red-role is like, and this is what no structural physics could have taught her (Proposition~\ref{p:persist}).

\subsection{Zombies}

A zombie is a duplicate of a conscious subject in respect of $P$ that lacks experience. By Theorem~\ref{t:modes} the zombie hypothesis is not ruled out a priori, since $P \to Q$ is not a priori; it is impossible, since $P \to Q$ is necessary. The conceivability of zombies has a precise source on the present account: in conceiving a zombie we hold fixed the structural mode of $P$ and suppose that the roles it describes are occupied by natures that do not ground experience. That supposition is coherent, because the structural mode does not say what occupies the roles; it is not a supposition about our world, because the roles in our world are occupied by natures that, in the arrangement in question, do ground experience. The zombie world is a world in which \emph{other} occupants play the roles. This is the Russellian loophole of MND-01-a, Assumption~4.1(c), viewed from the side of modes of presentation.

\subsection{The inverted spectrum}

MND-22 argues that spectrum inversion is possible but undetectable in principle, and therefore without consequence for functionalist theories of mind. The present account explains both halves. Inversion is possible because the intrinsic character of the occupant of a place in a quality space is not fixed by the place (MND-10-a, Proposition~6.1), and functional organization can be realized by arrangements whose intrinsic natures differ. It is undetectable because detection is by behavior and report, and behavior and report depend on the structural role of the arrangement, which inversion leaves unchanged. The two modes of presentation come apart here in the clearest way: two subjects can share everything that can be presented structurally and differ in what each is acquainted with. Note that inversion requires a difference in the occupants, not merely in how one occupant is presented; on the present account, an inverted subject in our world would be one whose qualia-grounding arrangements have different intrinsic natures, not one who thinks differently about the same natures.

\subsection{The meta-problem}

Why do we judge that there is a hard problem? The present account gives a two-part answer. First, our cognitive economy contains two systems of concepts for the same subject matter: theoretical concepts, acquired through structural description, and introspective concepts, acquired through acquaintance. That we have two such systems, and that they are not linked by a priori connections, is a structural fact about our cognitive architecture and is open to cognitive science. Second, when we reflect on the relation between the two systems we find that no derivation leads from the first to the second, and since in every other domain derivation is available (as in the water case), we judge that something special is going on. The judgment is correct: something special is going on, namely that only here does one of our systems of concepts present its subject matter intrinsically. The judgment is incorrect only if it is taken to show that the subject matter is special.

This answer satisfies the constraint of \citet{Chalmers2018} that a solution to the meta-problem should be connected to a solution to the hard problem. It does not reduce the hard problem to the meta-problem, as illusionism does \citep{Frankish2016}. The problem intuitions are explained by a difference in modes of presentation that is real and that involves the intrinsic character of experience, which illusionism denies. It is consistent with the result of MND-20, that introspection is reliable about current phenomenal states and unreliable about their causes: introspection delivers the intrinsic mode reliably and is silent about the structural roles of what it delivers, which is the asymmetry of Proposition~\ref{p:asym}.

\section{Objections and replies}\label{s:obj}

\begin{objection}[The master argument]\label{o:master}
Either the difference between the structural and intrinsic modes is physically explicable or it is not. If it is, a zombie, which shares everything physical, would have the same difference, and so the difference cannot explain our epistemic situation, which includes acquaintance with experience. If it is not, the account relocates the gap instead of explaining it.
\end{objection}

\begin{reply}
The objection, due to \citet{Chalmers2007}, assumes that \enquote{physically explicable} means explicable from structural truths. Take that sense. Then the difference between the modes is partly explicable and partly not. That we have two systems of concepts, that they are conceptually isolated, and that we deploy one in theory and the other in introspection are structural facts, and a zombie, which by stipulation shares the structural truths, would share them. What the zombie would not share is the relatum of acquaintance: in the zombie world the roles that ground our introspective concepts are occupied by natures that ground no experience, and so the zombie's introspective concepts present nothing intrinsically. The difference between us and the zombie is therefore in the occupants, and so is not explicable from structural truths. On the second horn, then, the account does not explain the gap from structural truths. But this is no objection to the account, since it says that the gap cannot be explained from structural truths, and why. What it does not do is posit a non-physical fact: the occupants are physical in the sense of MND-02, and the zombie world is not a physical duplicate of ours in that sense, since its occupants differ. The master argument has force against a view on which the referents of phenomenal concepts are fully characterized structurally; the present view is not of that kind.
\end{reply}

\begin{objection}[Revelation]\label{o:revelation}
If acquaintance reveals the nature of phenomenal properties, and physics reveals the structure of physical roles, then a subject possessing both could tell a priori whether a phenomenal property could occupy a given role. Either it could, and the gap closes, or it could not, and Russellian monism is false.
\end{objection}

\begin{reply}
The objection assumes that knowledge of what a property is in itself, together with knowledge of a role, yields a priori knowledge of whether the property occupies the role. It does not. That a given property occupies a given role is a contingent matter of which property is lawfully connected with which, and knowing two things fully does not tell one how they are related in the world: knowing the intrinsic character of an object and the description of a role does not tell one whether the object plays it. \citet{Goff2017} uses a doctrine of revelation against type-B physicalism, arguing that phenomenal concepts reveal the essence of their referents, and so cannot be concepts of properties whose essence is neural and functional. The present account is consistent with a partial revelation: acquaintance reveals some of what an experience is in itself, not that it does not occupy a role. It is also consistent with Goff's conclusion against type-B physicalism, since the account is not type-B physicalism. The phenomenal property is not identified with a structurally characterized property; its ground is the occupant of a structural role.
\end{reply}

\begin{objection}[Grounding is not identity]\label{o:ground}
On the protophenomenal version of Russellian monism, the intrinsic natures that occupy physical roles are not phenomenal. Experience presents phenomenal properties, which are grounded in those natures but are not identical to them. So experience does not present the same thing that physics describes, and the slogan \enquote{two modes of presentation of the same facts} is false of the view.
\end{objection}

\begin{reply}
The objection is right that experience does not present the protophenomenal natures of individual microphysical occupants, and it is right that phenomenal properties are grounded in them rather than identical to them. But the claim of MND-01 is not that one property is presented twice. It is that one portion of reality is presented twice, in the sense defined in Section~\ref{s:frame}: what $Q$ concerns is wholly grounded in what $P$, read as concerning the actual occupants, concerns, and adds nothing to it (Lemma~\ref{l:same}). Physics presents the arrangement as a structure of role-occupants; experience presents what that arrangement is like, taken as a whole. There are two modes because the arrangement is presented once by its relations and once by its character. The individual natures in it are presented by neither mode as such, which is why the palette problem arises (MND-01-d, Section~2). That is a limitation of our knowledge of the grounding, not a second subject matter.
\end{reply}

\begin{objection}[Illusionism]\label{o:illusion}
\ref{p:acq} assumes that there is something with which we are acquainted, an intrinsic phenomenal character. The illusionist denies this: there is only the introspective representation of such a character, and the gap is a product of that representation's misdescribing its object.
\end{objection}

\begin{reply}
The present manuscript does not argue against illusionism; MND-10-a does, by a dilemma about the seeming that illusionism posits (MND-10-a, Lemma~4.2). But two points are worth making. First, the illusionist and the present account agree that the meta-problem has a structural solution, the existence of two isolated systems of concepts. They disagree about whether one of the systems presents anything intrinsically. Second, if the illusionist were right, \ref{p:acq} would fail, and Lemma~\ref{l:gap} would lose its proof, since $Q$ would then be a structural truth about representations. The gap would have to close, at least in principle, as the gap between water and H$_2$O closes. The persistence of the gap across every refinement of cognitive science is therefore evidence for \ref{p:acq}, and the illusionist owes an account of why a gap between two structural descriptions should be unclosable in principle.
\end{reply}

\begin{objection}[Ignorance]\label{o:ignorance}
The account is a version of the ignorance hypothesis of \citet{Stoljar2006}, with a specific conjecture about what we are ignorant of. But the ignorance hypothesis predicts that if we came to know the missing truths, the gap would close. The account says the gap cannot close. It is therefore inconsistent with the hypothesis it claims to specify.
\end{objection}

\begin{reply}
The account agrees with the ignorance hypothesis that $P$ omits a type of truth about the physical world, truths about the intrinsic natures of role-occupants. It disagrees on how those truths are known. For Stoljar, they are truths of a kind that a future science might discover and state. On the present account, the only access to them is acquaintance, and acquaintance presents only what the arrangements grounding our own experiences are like. The gap therefore closes for a subject, in the sense of Proposition~\ref{p:occ}, once she adds the occupant identifications to $P$; but those identifications are not available from the structural mode, and so the gap between the truths of physics, as stated, and the truths of experience, as known, never closes (Proposition~\ref{p:persist}). This is not inconsistent with the ignorance hypothesis but a specification of it on which the ignorance is irremediable by physics and partly remedied by acquaintance.
\end{reply}

\section{What the account does not establish}\label{s:not}

The account explains the gap given Russellian monism; it does not argue for Russellian monism. That argument is given in MND-01-a, from the two-dimensional argument and materialism, and in MND-01-c by a route that avoids the two-dimensional premises. Without Assumption~\ref{a:rm}, Lemma~\ref{l:gap} still holds, since its proof uses only \ref{p:struct-a}--\ref{p:acq}, but Theorem~\ref{t:modes} does not: a dualist can accept that $P \to Q$ is not a priori and deny that it is necessary.

The account does not say which arrangements ground which experiences. Proposition~\ref{p:occ} locates the a posteriori element in occupant identifications and Corollary~\ref{c:noresid} locates the remaining explanatory work in the grounding principles. Stating those principles is the task of the boundary and palette problems, which are problems of articulation rather than of possibility (MND-01-d, Proposition~2.1), and which Volume~IV of MND-01 treats in full.

The account does not make the intrinsic natures of fundamental entities phenomenal. Experience presents the intrinsic character of arrangements that ground subjects, and nothing in the account implies that there is anything it is like to be an electron. MND-18 shows that a view on which there were would face an unsolvable combination problem (MND-18-a, Theorem~3.2).

Finally, the account does not deny that the phenomenal concept strategy captures something true. Phenomenal concepts are recognitional and conceptually isolated, as MND-04 and MND-05 hold. The present account explains why they are isolated: they present their referents in the intrinsic mode, and \ref{p:closure} forbids an a priori bridge from the structural mode to the intrinsic one.

\section{Conclusion}\label{s:concl}

The explanatory gap is the a priori non-derivability of phenomenal truths from physical truths. It follows from the fact that physical truths present their subject matter structurally and phenomenal truths, as known to the subject, present it intrinsically, together with the closure of a priori reasoning under structural content. Given the Russellian monism of MND-01, the two kinds of truth concern the same portion of reality, so the gap is wholly epistemic; it persists under every structural extension of physics; and its a posteriori residue consists of occupant identifications, which pair a role with what occupies it and so do not invite the question why a property characterized one way should also be characterized another. The gap is a gap between modes of presentation.

\begin{thebibliography}{99}\small
\bibitem[Balog(2012)]{Balog2012} Balog, K. (2012). Acquaintance and the mind--body problem. In C.~Hill and S.~Gozzano (Eds.), \emph{New Perspectives on Type Identity}. Cambridge University Press.
\bibitem[Block and Stalnaker(1999)]{BlockStalnaker1999} Block, N. and Stalnaker, R. (1999). Conceptual analysis, dualism, and the explanatory gap. \emph{The Philosophical Review}, 108(1), 1--46.
\bibitem[Chalmers(1996)]{Chalmers1996} Chalmers, D.~J. (1996). \emph{The Conscious Mind: In Search of a Fundamental Theory}. Oxford University Press.
\bibitem[Chalmers(2007)]{Chalmers2007} Chalmers, D.~J. (2007). Phenomenal concepts and the explanatory gap. In T.~Alter and S.~Walter (Eds.), \emph{Phenomenal Concepts and Phenomenal Knowledge} (pp.~167--194). Oxford University Press.
\bibitem[Chalmers(2018)]{Chalmers2018} Chalmers, D.~J. (2018). The meta-problem of consciousness. \emph{Journal of Consciousness Studies}, 25(9--10), 6--61.
\bibitem[Chalmers and Jackson(2001)]{ChalmersJackson2001} Chalmers, D.~J. and Jackson, F. (2001). Conceptual analysis and reductive explanation. \emph{The Philosophical Review}, 110(3), 315--360.
\bibitem[Frankish(2016)]{Frankish2016} Frankish, K. (2016). Illusionism as a theory of consciousness. \emph{Journal of Consciousness Studies}, 23(11--12), 11--39.
\bibitem[Frege(1892)]{Frege1892} Frege, G. (1892). \"Uber Sinn und Bedeutung. \emph{Zeitschrift f\"ur Philosophie und philosophische Kritik}, 100, 25--50.
\bibitem[Goff(2017)]{Goff2017} Goff, P. (2017). \emph{Consciousness and Fundamental Reality}. Oxford University Press.
\bibitem[Levine(1983)]{Levine1983} Levine, J. (1983). Materialism and qualia: the explanatory gap. \emph{Pacific Philosophical Quarterly}, 64(4), 354--361.
\bibitem[Levine(1993)]{Levine1993} Levine, J. (1993). On leaving out what it's like. In M.~Davies and G.~W. Humphreys (Eds.), \emph{Consciousness: Psychological and Philosophical Essays}. Blackwell.
\bibitem[Levine(2001)]{Levine2001} Levine, J. (2001). \emph{Purple Haze: The Puzzle of Consciousness}. Oxford University Press.
\bibitem[Lewis(1970)]{Lewis1970} Lewis, D. (1970). How to define theoretical terms. \emph{The Journal of Philosophy}, 67(13), 427--446.
\bibitem[Loar(1990)]{Loar1990} Loar, B. (1990). Phenomenal states. \emph{Philosophical Perspectives}, 4, 81--108.
\bibitem[Papineau(2002)]{Papineau2002} Papineau, D. (2002). \emph{Thinking about Consciousness}. Oxford University Press.
\bibitem[Russell(1911)]{Russell1911} Russell, B. (1911). Knowledge by acquaintance and knowledge by description. \emph{Proceedings of the Aristotelian Society}, 11, 108--128.
\bibitem[Russell(1927)]{Russell1927} Russell, B. (1927). \emph{The Analysis of Matter}. Kegan Paul.
\bibitem[Stoljar(2005)]{Stoljar2005} Stoljar, D. (2005). Physicalism and phenomenal concepts. \emph{Mind and Language}, 20(5), 469--494.
\bibitem[Stoljar(2006)]{Stoljar2006} Stoljar, D. (2006). \emph{Ignorance and Imagination: The Epistemic Origin of the Problem of Consciousness}. Oxford University Press.
\bibitem[Tye(2009)]{Tye2009} Tye, M. (2009). \emph{Consciousness Revisited: Materialism without Phenomenal Concepts}. MIT Press.
\end{thebibliography}

\end{document}
